\documentclass[10pt,a4paper]{amsart}
\usepackage[margin=19mm]{geometry}
\usepackage{amsmath,amssymb,mathtools}
\usepackage[scr=rsfs,cal=euler]{mathalfa}
\usepackage{xfrac}
\usepackage[unicode,hidelinks,bookmarks=false]{hyperref}
\newcommand{\ie}{i.e.\ }
\newcommand{\cf}{cf.\ }
\newcommand{\resp}{respectively\ }
\newcommand{\Subsection}{\subsection}
\providecommand{\nobreakdash}{\nobreak\hbox{-}}
\setlength{\emergencystretch}{2em}
\multlinegap0pt
\usepackage{bm}
\usepackage{bbm}
\usepackage{dsfont}
\usepackage{xspace}
\usepackage{cancel}
\usepackage{mathtools}
\usepackage{amsthm}
\makeatletter
\@ifundefined{theorem}{\newtheorem{theorem}{Theorem}}{}
\@ifundefined{lemma}{\newtheorem{lemma}[theorem]{Lemma}}{}
\@ifundefined{proposition}{\newtheorem{proposition}[theorem]{Proposition}}{}
\makeatother
\newcommand\E{\mathbb{E}}
\newcommand\R{\mathbb{R}}
\newcommand\Tr{{\operatorname{Tr}}}
\newcommand\bg{\bm{g}}
\newcommand\bu{\bm{u}}
\newcommand\bw{\bm{w}}
\newcommand\eps{\varepsilon}
\newcommand\hg{\hat{\bm g}}
\newcommand\hu{\hat{\bm u}}
\newcommand\proj{{\operatorname{P}}}
\newcommand\rem{\tilde{O}\left(n^{-3/2}\right)}
\newcommand{\Cov}{\mathrm{Cov}}
\newcommand{\Var}{\mathrm{Var}}
\title[Nodal Count for Orthogonally Invariant Ensembles]{Nodal Count for Orthogonally Invariant Ensembles}
\author{Lior Alon, Dan Mikulincer,, John Urschel}
\date{Source: arXiv:2511.02784v1 (CC BY 4.0)}
\begin{document}
\maketitle

\noindent\emph{Source and licence.} Nodal Count for Orthogonally Invariant Ensembles. Version arXiv:2511.02784v1 (CC BY 4.0), distributed under the Creative Commons Attribution 4.0 International licence, as recorded in the arXiv licence metadata for this exact version and shown on the abstract page https://arxiv.org/abs/2511.02784v1. Full rights notice: licenses/arxiv-2511.02784-CC-BY-4.0.txt.\par\smallskip

\noindent\textit{Editorial scope.} Counted unit: the Proposition whose source label is \texttt{prop:main\_prop} in the article (source0.tex lines 714--724, source label \texttt{prop:main\_prop}), delivered together with the article's own complete original proof of it (source0.tex lines 729--820, one \texttt{proof} environment of 92 lines). No mathematics is written, repaired or re-derived: statement and proof are the article's own text line for line. The proof is detached from the statement in the source: the article states the result at lines 714--724 and prints its proof at lines 729--820, and the proof environment's own header names the statement, so the association is the article's own. Required context retained uncounted: lm:typical (lines 512--516); lem: signArcsine (lines 703--710). Every definition, display and label of those lines is retained unchanged. Omitted from this package and not redistributed: 1--511, 517--702, 711--713, 725--728, 821--1046. Nothing inside a retained excerpt is omitted or altered: every retained body is delivered exactly as the article's lines are written, so no omission receipt of any kind accompanies this package. Citations: the retained excerpts carry no citation command at all, so no bibliography entry of the article is transcribed and no reference is redistributed. Reference adaptation: none was needed and none was made; every reference inside the delivered unit resolves inside it, so no unresolved reference or citation is produced. Printed slips kept literal: no printed word, symbol, hypothesis or number of the counted statement or of its proof was altered. Layout and class: the article's own class is \texttt{amsart} with packages including amsaddr, inputenc, svg, hyperref, url, amssymb, amscd, amsthm, setspace, enumerate, graphicx, mathtools, tcolorbox, subcaption; the wrapper is the lane's pinned \texttt{amsart} editorial wrapper at 10pt on A4, which restates the theorem-like environments the retained text needs and adds the article's own macro definitions that the retained text needs (\texttt{\textbackslash Cov}, \texttt{\textbackslash E}, \texttt{\textbackslash R}, \texttt{\textbackslash Tr}, \texttt{\textbackslash Var}, \texttt{\textbackslash bg}, \texttt{\textbackslash bu}, \texttt{\textbackslash bw}, \texttt{\textbackslash eps}, \texttt{\textbackslash hg}, \texttt{\textbackslash hu}, \texttt{\textbackslash proj}, \texttt{\textbackslash rem}), with extra wrapper packages (bm, bbm, dsfont, xspace, cancel, mathtools). The delivered numbering is the unit's own. Assets: no retained span contains a figure environment, a graphics-inclusion command, a PDF-page inclusion, a table or a TikZ picture; no image file is copied into this package, the package declares no asset dependency and no image file is redistributed with it. Licence: version arXiv:2511.02784v1 of this article is distributed under the Creative Commons Attribution 4.0 International licence, as recorded in the arXiv licence metadata for this exact version and shown on the abstract page https://arxiv.org/abs/2511.02784v1. A full rights notice is delivered at licenses/arxiv-2511.02784-CC-BY-4.0.txt. Characters: every retained excerpt is pure ASCII, so the delivered engine, XeLaTeX with its default Latin Modern text font, reports no missing character.
\begin{lemma}\label{lm:typical}
Let $\bm{x} \in \mathbb{R}^{n}$ be a fixed vector with $\|\bm{x}\|_2 \le n^{1/2} \log^{c_1} n$, $\|\bm{x}\|_4 \le n^{1/4} \log^{c_1} n$ and $G \in \mathbb{R}^{n \times p}$ be a standard Gaussian matrix, for some fixed constants $c_1$ and $p$. Then there exists a constant $c_2$ such that $\|G\|_{\max} \le \log^{c_2} n$ and
\[\left\| G^T \mathrm{diag}(\bm{x})^k G - \E\big[ G^T \mathrm{diag}(\bm{x})^k G \big] \right\|_{\max} \le \sqrt{n} \log^{c_2} n\]
for $k = 0,1,2$ with probability $1 - n^{-\log n}$.
\end{lemma}

\begin{lemma}\label{lem: signArcsine}
Let $(X,Y)\in \R^2$ be a centered Gaussian vector with unit variances and $\E[XY] = \rho$.
Then
\[
\E\!\big[\mathrm{sgn}(X)\,\mathrm{sgn}(Y)\big]
=\frac{2}{\pi}\arcsin(\rho)=\frac{2}{\pi}\left(\rho+\frac{\rho^3}{6}+\frac{3\rho^5}{40}+O(\rho^7)\right).
\]
\end{lemma}

\begin{proposition}\label{prop:main_prop}
Let $B$ be a real $n\times n$ matrix and let $P$ be a rank-one orthogonal projection.
Assume that $\|(\mathrm{I}-P)B(\mathrm{I}-P)\|_{F}^2=n+\tilde{O}(1)$ and that the operator norms
$\|B\|$, and $\|B^{T}P\|$ are $\tilde{O}(1)$. Further assume $\Tr(B)=\tilde{O}(1)$.
Let $\bu,\hu$ denote the first two columns of a Haar-random orthogonal matrix in $\mathcal{O}(n)$. Then
\[
\E\!\left[\mathrm{sgn}\!\big(\langle P\bu,\hu\rangle\,
\langle B\bu,\hu\rangle\big)\right]
=\frac{2^{3/2}}{\pi^{3/2}}\Tr(PB)n^{-1/2}+\tilde{O}(n^{-3/2}).
\]
\end{proposition}

\begin{proof}[Proof of Proposition \ref{prop:main_prop}]
Our proof is conducted in several steps. We first prove an analog result for quadratic forms in Gaussians. We then explain how to reduce, through conditioning, the computation on the orthogonal group to Gaussian space. We then estimate the error by integrating over the conditioned variables.

\emph{Step 0: Introduce the Gaussian setting.} Let $\bg$ and $\hg$ be two independent standard Gaussian vectors in $\R^n$. By rotational invariance of $(\bg,\hg)$, we may and will assume $\proj=e_{1}e_{1}^T$ is the projection onto the first coordinate. 
The well-known construction of the first two Haar columns is
\[
(\bu,\hu) \;\stackrel{\mathrm{law}}=\; \left(\frac{\bg}{\|\bg\|}, \frac{Q\,\hg}{\|Q\,\hg\|}\right),
\quad
Q=Q_{\bg}:=\mathrm{I}_n-\frac{\bg\bg^\top}{\|\bg\|^2}.
\]
The map
$(u,v)\mapsto \mathrm{sgn}\big(\langle Pu,v\rangle\,\langle Bu,v\rangle\big)$
is invariant under \emph{independent scalings} of $u$ and $v$.
Hence
\[
\mathrm{sgn}\!\big(\langle P\bu,\hu\rangle\,\langle B\bu,\hu\rangle\big)
\ \stackrel{\mathrm{law}}=\ 
\mathrm{sgn}\!\big(\langle P\bg,Q\hg\rangle\,\langle B\bg,Q\hg\rangle\big)=\ 
\mathrm{sgn}\!\big(\, g_{1}\,\langle Qe_1,\hg\rangle\,\langle QB\bg,\hg\rangle\big).
\]

\emph{Step 1: Condition on $\bg$ and reduce to two linear forms of $\hg$.}
Fix $\bg$ and set
\[
u:=\mathrm{sgn}(g_1)Qe_1,\qquad v:=QB\bg, \qquad
X:=\frac{\langle u,\hg\rangle}{\|u\|},\qquad
Y:=\frac{\langle v,\hg\rangle}{\|v\|}.
\]
For almost any value of $\bg$, $(X,Y)\mid\,\bg$ is bivariate normal, centered, with
\[
\Var(X\mid \bg)=1,\quad \Var(Y\mid \bg)=1,\quad
\rho:=\Cov(X,Y\mid \bg)=\frac{\langle u,v\rangle}{\|u\|\|v\|}=\mathrm{sgn}(g_1)\frac{\langle Qe_1,QB\bg\rangle}{\|Qe_1\|\|QB\bg\|},
\]
so that, by Lemma Lemma \ref{lem: signArcsine},
\[\E_{\hg}\big(\mathrm{sgn}\!\big(\langle P\bu,\hu\rangle\,\langle B\bu,\hu\rangle\big)\big)=\E_{\hg}\big(\mathrm{sgn}(X)\mathrm{sgn}(Y)\big)=\frac{2}{\pi}\left(\rho+\frac{\rho^3}{6}\right)+O(\rho^5).\]
It is now a matter of bounding the random correlation \(\rho\). Notice that $\langle Qe_1,QB\bg\rangle=\langle Qe_1,B\bg\rangle$ due to $Q^2=Q=Q^T$. The fact that $\|Qx\|^2=\|x\|^2-\frac{|\langle x,\bg\rangle|^2}{\|\bg\|^2}$ for any $x\in\R^n$ allows to write
\[
\mathrm{sgn}(g_{1})\rho
=\frac{\|\bg\|^2\langle Qe_1,B\bg\rangle}{\sqrt{(\|\bg\|^2-g_{1}^2)(\|\bg\|^2\|B\bg\|^2-|\langle \bg,B\bg\rangle|^2)}}
=\frac{\|\bg\|^2(B\bg)_1-g_{1}\langle \bg,B\bg\rangle}{\sqrt{(\|\bg\|^2-g_{1}^2)(\|\bg\|^2\|B\bg\|^2-|\langle \bg,B\bg\rangle|^2)}}.
\]
Define the fluctuations 
\begin{equation} \label{eq:errors}
   \eps:=\frac{\|\bg\|^2-n}{n},\quad \delta:=\frac{\|B\bg\|^2-n}{n},\quad \eta:=\frac{\langle\bg,B\bg\rangle}{n}.
\end{equation}
Notice that $\E[\eps]=0,\ \E[\delta]=n^{-1}(\|B\|_{F}^2-n),$ and $\E[\eta]=n^{-1}\Tr(B)$. The assumptions that $\|B\|_{F}^2-n=O(1)$ and $\Tr(B)=\tilde{O}(1)$, together with the Gaussian concentration from Lemma \ref{lm:typical}, show that with high probability (e.g. $1-n^{-\log n}$) the random variables $\eps,\delta,\eta$ are $\tilde{O}(n^{-1/2})$. Define the Gaussian 
\[Z=(B\bg)_1=\langle B^Te_1,\bg\rangle.\]
Notice that $\|B^Te_1\|=\|B^TPe_1\|\le\|B^TP\|=\tilde{O}(1)$ by assumption. So $g_1$ and $Z$ are centered Gaussians with variances $1$ and $\tilde{O}(1)$ respectively. In particular, with high probability $g_1$ and $Z$ are $\tilde{O}(1)$. We may conclude that with high probability (e.g. $1-n^{-\log n}$),
\begin{align*}
\mathrm{sgn}(g_{1})\rho=& \frac{n\,(1+\eps)Z-n\,\eta\, g_1}{\sqrt{\big(n\,(1+\eps)-g_1^2\big)\big(n^2(1+\eps)(1+\delta)-n^2\eta^2\big)}}\\
= &    n^{-1/2}\frac{Z-\eta g_1+\rem}{\sqrt{\big(1+\eps+g_{1}^2\,n^{-1}\big)\big(1+\eps+\delta+\eps\,\delta-\eta^2)\big)}}\\
= &    n^{-1/2}\frac{Z-\eta g_1+\rem}{\sqrt{1+2\eps+\delta+g_{1}^2\,n^{-1}+\eps\,\delta-\eta^2+\rem}}\\
= &    n^{-1/2}\big(Z-\eta g_1\big)\big(1-\frac{1}{2}(2\eps+\delta+g_{1}^2\,n^{-1}+\eps\,\delta-\eta^2)+\frac{3}{8}(2\eps+\delta)^2+\rem\big)\\
\\
= &    n^{-1/2}\big[Z\big(1-\eps-\frac{\delta}{2}-\frac{1}{2}(g_{1}^2\,n^{-1}+4\eps\,\delta-\eta^2+3\eps^2+\delta^{2}/4)\big)-\eta g_{1}\big(1-\eps-\frac{\delta}{2}\big)\big]+\tilde{O}(n^{-2})\\
= &    n^{-1/2}Z-n^{-1/2}\big(Z(\eps+\frac{\delta}{2})+\eta g_{1}\big)+W+\tilde{O}(n^{-2}),
\end{align*}
where $W$ is with high probability $\rem$
\[W:=n^{-1/2}\big(\eta g_1(\eps+\frac{\delta}{2})-Z\frac{1}{2}(g_{1}^2\,n^{-1}+4\eps\,\delta-\eta^2+3\eps^2+\delta^{2}/4)\big).\]
In particular, with probability at least  $\ 1-n^{-\log n}$,
\[\rho=\tilde{O}(n^{-1/2}),\quad\text{and} \quad\rho=n^{-1/2}\mathrm{sgn}(g_{1})\left(Z-Z\eps-Z\frac{\delta}{2}+\eta g_1\right)+\rem. \]
Since $|\mathrm{sgn}| \leq 1$, we can now calculate up to $\rem$,
\begin{align*}
    \E\!\left[\mathrm{sgn}\!\big(\langle P\bu,\hu\rangle\,
\langle B\bu,\hu\rangle\big)\right]
= & \frac{2}{\pi}\E_{\bg}[\rho]+\rem.
\end{align*}
\emph{Step 2: Integration over $\bw=(\mathrm{I}_{n}-e_{1}e_{1}^{T})\bg$.}
To do so, we let $\bw=(0,g_{2},\ldots,g_{n})$, so that its $n-1$ non-zero entries form a standard $n-1$ dimensional Gaussian independent of $g_1$, and we first integrate over $\bw$. We can write 
\[\eps=\frac{\|\bw\|^2+g_{1}^2-n}{n},\quad \delta:=\frac{\|B\bw\|^2+g_{1}B_{11}-n}{n},\quad \eta:=\frac{B_{11}g_{1}^2+\langle\bw,B\bw\rangle+\langle\bw,Be_{1}\rangle+\langle e_{1},B\bw\rangle}{n},\]
and
\[Z=B_{11}g_{1}+\langle B^{T}e_{1},\bw\rangle.\]
Notice that by our assumptions, $\E_{\bw}\|B\bw\|^2=\|(\mathrm{I}-P)B(\mathrm{I}-P)\|_{F}^2=n+\tilde{O}(1)$ and $\E_{\bw}\langle \bw,B\bw\rangle=\Tr((\mathrm{I}-P)B(\mathrm{I}-P))=\Tr(B)-\Tr(PB)=\tilde{O}(1)$. Using these observations and the symmetry $\bw\mapsto-\bw$, we get  
\begin{align*}
\E_{\bw}[Z]=& B_{11}g_{1},\\
\E_{\bw}[Z\eps]=& \E_{\bw}[B_{11}g_{1}\eps]=\frac{B_{11}g_{1}(g_{1}^2-1)}{n}=\tilde{O}(1/n),\\
\E_{\bw}[Z\delta]=& \E_{\bw}[B_{11}g_{1}\delta]=\frac{B_{11}g_{1}(B_{11}g_{1}+\tilde{O}(1))}{n}=\tilde{O}(1/n),\\
\E_{\bw}[\eta]=& \frac{B_{11}g_{1}^2+\tilde{O}(1)}{n}=\tilde{O}(1/n).
\end{align*}
We conclude that,
\begin{align*}
\E_{\bg}[\rho]=\E_{g_{1}}\left[\E_{\bw}(\rho)\right]= & n^{-1/2}\Tr(PB)\E_{g_{1}}\left[|g_{1}|\right]+\rem\\ = & \sqrt{\frac{2}{\pi}}n^{-1/2}\Tr(PB)+\rem,  
\end{align*}
and therefore
\begin{equation*}
 \E\!\left[\mathrm{sgn}\!\big(\langle P\bu,\hu\rangle\,
\langle B\bu,\hu\rangle\big)\right]
=  \frac{2}{\pi}\E_{\bg}[\rho]+\rem
= \frac{2^{3/2}}{\pi^{3/2}}\Tr(PB)n^{-1/2}+\rem. 
\end{equation*}
This proves the claim.
\end{proof}

\end{document}
